\documentclass[10pt]{sigplanconf}

% The following \documentclass options may be useful:

% preprint      Remove this option only once the paper is in final form.
% 10pt          To set in 10-point type instead of 9-point.
% 11pt          To set in 11-point type instead of 9-point.
% authoryear    To obtain author/year citation style instead of numeric.

\usepackage{amsmath}
\usepackage{listings}
\usepackage{xcolor}
\usepackage{hyperref}

\lstdefinelanguage[80]{Smalltalk}{
  morekeywords={true,false,self,super,nil,thisContext},
  sensitive=true,
  morecomment=[s]{"}{"},
  morestring=[d]',
  style=SmalltalkStyle
}
\lstdefinestyle{SmalltalkStyle}{
  literate=%
  % {:=}{{$\gets$}}1
  {^}{{$\uparrow$}}1
  % {…}{{\coderem}}1
  {…}{{\ldots}}1
}

\lstset{
  defaultdialect=[80]{Smalltalk},
  language=Smalltalk,
  basicstyle=\footnotesize\ttfamily,%
  numbers=left,%
  numbersep=5pt,%
  numberstyle=\tiny\color{gray},%
  emphstyle=\bfseries,%
  breakatwhitespace=false,%
  showstringspaces=false,%
  columns=flexible,%
  frame=b,%
  rulecolor=\color{lightgray} }


\begin{document}

\special{papersize=8.5in,11in}
\setlength{\pdfpageheight}{\paperheight}
\setlength{\pdfpagewidth}{\paperwidth}

\conferenceinfo{ICOOOLPS'15}{July 6, 2015, Prague, Czech Republic} 
\copyrightyear{2015} 
\copyrightdata{978-1-nnnn-nnnn-n/yy/mm} 
\doi{nnnnnnn.nnnnnnn}

\exclusivelicense                % ACM gets exclusive license to publish, 
                                 % you retain copyright

\preprintfooter{Preprint}   % 'preprint' option specified.

\title{Algorithmic VM Primitives in a High-Level Language}
\subtitle{Applying a Tracing JIT to Minimize the Need for Primitives}

\authorinfo{Tim Felgentreff\and Tobias Pape\and Robert Hirschfeld}
           {Hasso Plattner Institute, University of Potsdam}
           {\{firstname.lastname\}@hpi.uni-potsdam.de}

\maketitle

\begin{abstract}
  When implementing virtual machines, besides the interpreter and
  optimization facilities, we have to include a set of primitive
  functions that the client language can use. Some of these implement
  truly primitive behavior, such as arithmetic operations. Other
  primitive functions, which we call \emph{algorithmic primitives},
  are expressible in the client language, but are implemented in the VM
  to improve performance.

  However, having many primitives in the VM makes it harder to
  maintain them, or re-use them in alternative VM implementations for
  the same language. With the advent of efficient tracing just-in-time
  compilers we believe the need for algorithmic primitives to be much
  diminished, allowing more of them to be implemented in the client
  language.

  In this work, we investigate the trade-offs when creating
  primitives, and in particular how large a difference remains between
  primitive and hosted function run times in VMs with tracing
  just-in-time compiler. To that end, we extended the R/SqueakVM, a VM
  for Squeak/Smalltalk written in RPython. We compare primitive
  implementations in C, RPython, and Smalltalk, showing that due to
  the tracing JIT the performance gap can be significantly reduced.
\end{abstract}

\category{CR-number}{subcategory}{third-level}

% general terms are not compulsory anymore, 
% you may leave them out
\terms
term1, term2

\keywords
keyword1, keyword2

\section{Introduction}
Dynamically typed programming languages are usually implemented as
process Virtual Machines (VMs). With the increasing popularity of
these languages, VM implementation techniques have evolved to the
point where most mature VMs utilize some form of dynamic compilation,
such as a method just-in-time compiler (JIT), or a tracing
JIT. Another large part of a VM are basic behavioral blocks, so called
primitives. Some of these primitives are unavoidable, because their
behavior cannot adequately be expressed in the implemented client
language. But some primitives, which we call \emph{algorithmic
  primitives}, exist only because, supposedly, primitive functions
evaluate faster than identical functions implemented in the client
language. We discuss this in more detail in
\autoref{sec:prim-build-blocks}.

There are few languages which have just one implementation. Some of
those language implementations are not (only) source compatible, but
for example share a bytecode set or a standard library. As a result,
every primitive in the reference implementation has to be mirrored in
all the other implementations. Furthermore, the primitive
implementations cannot usually be debugged from the high-level
language --- if a bug is suspected in a primitive implementation, the
debugging process has to cross language boundaries and thus becomes
more cumbersome. Thus, we argue, it is desirable to reduce the number
of primitives to ease VM development and maintenance.

In this work we investigate the trade-offs around implementing
behavior in the VM versus the client language, and make the following
contributions to this field:
\begin{itemize}
\item We discuss possible trade-offs when implementing behavior as
  primitive in contrast to in the hosted language in
  \autoref{sec:trade-offs-between}.
\item We quantify the performance gap between functionality
  implemented as primitive in contrast to in the hosted language in
  \autoref{sec:measurements}.
\end{itemize}

\section{Primitive Building Blocks in Virtual Machines}
\label{sec:prim-build-blocks}
When invoking behavior in languages implemented as VM, there are ways
to call functions which are not implemented in the client
language. Some of those functions are grouped into libraries and
called using a Foreign Function Interface (FFI). Other functions,
commonly called primitives, are a fixed part of the VM.

The most common primitive behaviors are often available as
bytecodes. This reduces the number of function names to be saved near
the call site, and may elide the lookup of the function name, thus
saving space and time. Depending on the language implementation, such
bytecodes might be the only way to achieve some behavior, e.g. integer
addition on the Java Virtual Machine (JVM).

In most languages, there is no difference between calling a primitive
or calling a hosted function, i.e. primitive invocation is
transparent. In this work, we use R/SqueakVM, an implementation of
Squeak/Smalltalk in RPython. RPython is the language and toolchain
that also powers the PyPy Python implementation

... todo describe RPython in 1--2 sentences

The original Squeak VM is written in a static subset of Smalltalk
called \emph{Slang}. This code is statically translated into C and
compiled to create the virtual machine. In addition, many methods in
Squeak that invoke primitive behavior have so-called \emph{fallback
  code} that is executed when the primitive fails. Consider the
following method:
\begin{lstlisting}
bitXor: n 
  | norm |
  <primitive: 'primDigitBitXor' module:'LargeIntegers'>
  norm := n normalize.
  ^ self
      digitLogic: norm
      op: #bitXor:
      length: (self digitLength max: norm digitLength)
\end{lstlisting}

In Squeak/Smalltalk, a method can have an annotation, called pragma
(line 3). The pragma can indicate a primitive by number or name. In
the above example, the primitive behavior is implemented as
\verb!primDigitBitXor! in the \verb!LargeIntegers! module. When the
interpreter encounters a method with such a pragma, it dispatches to
the primitive instead of running the method source. Only if the
primitive fails, either because it is not implemented in the executing
VM, or because the primitive codes assumptions about its arguments
fail, the VM invokes the method’s body. In some cases the body may
just generate a run-time exception. In the above snippet, however, the
code implements the same behavior is the primitive.

In file-organized languages, such as Python, calling primitives is
like calling an imported function. The only difference is that the
imported function is not implemented in the hosted
language. Primitives in such languages are usually documented in text
form, explaining a functions signature, arguments, side-effects, and
return value.

In Self, primitives are indicated by their selectors (slot
names). Every access to a slot whose name starts with underscore (\_)
should call a primitive function. If that function fails, the VM
raises an error to be caught by the caller.

In the following paragraphs, we present a grouping of primitives based
on their behavior, ordered according to the necessity for having them
in a VM.

\paragraph{Memory Primitives}
\label{sec:memory-primitives}
This set of primitives is closely related to the memory design of the
programming language. In object-oriented languages, memory primitives
provide means to create new objects, access their fields, their size,
and trigger garbage collection. The behavior supplied by these
primitives might be implemented as bytecodes, because they are used
frequently. 

\paragraph{Mathematical Primitives}
While arithmetics and boolean operations can be implemented using
memory access, branching, and constants, they are used often and have
a magnitudes faster implementation in the actual machine.  This group
of primitives consists of basic arithmetic operations such as
addition, multiplication, and inverses, but also trigonometric
functions, modulo, truncation, exponential functions, numeric type
conversion, and root functions. The binary value primitives are and,
or, xor, bit-shift, sometimes also bit-invert. Another bigger subgroup
of mathematical primitives are comparison operations, such as smaller,
smaller or equal, equality, identity, and their inverses.

\paragraph{Foreign Function Interface (FFI)}
Usually, VMs offer a way to interface with existing libraries. The
greatest common denominator among FFIs are C call- ing conventions. A
FFI usually requires explicit naming of the library to find the
functions in, and marshalling for function arguments and demarshalling
of return values. While in primitives, the callee is responsible for
type conversions, in FFI, the caller has to explicitly convert types.

\paragraph{System Calls}
Most process VMs offer a way to access system resources, such as the
file system, system clock, or the network. Some offer an abstracted
API trying to merge similar concepts of different operating systems,
giving access to the greatest common denominator of
functionality. This way, they try to promote operating system
independent programming. The goal of a programming language
determines, whether this set of primitives includes input and output
(IO)-operations, or not.

\paragraph{VM Specific Primitives}
Depending on the implementation, the VM might offer additional entry
points, for example to register callbacks for mouse and keyboard
handling, for green threading, timed behavior, or reflection on VM
internals such as garbage collection information.

\paragraph{Algorithmic Primitives}
Some behavior can be sufficiently expressed in the hosted language,
but is never the less available as primitive. Primitives are generally
considered to be faster than the same function implemented in the
hosted language. This thesis focuses on this group of primitives,
showing why current optimization techniques are well suited and in how
far the ex- ecution time gap has closed as a result. Reasons for
implementing behavior as primitive, or in the hosted language, are
explored in the following chapter.  Common examples are copying blocks
of memory from one collection to an- other, simple loops which are run
regularly, sorting algorithms, and integer arithmetics for arbitrarily
large numbers.

\section{Trade-Offs Between Implementation Locations}
\label{sec:trade-offs-between}

Whenever programmers make the decision to implement a function as a
primitive, they have to weigh the behavior, speed, and security
against changeability, reuse, and ease of implementation.

\subsection{Primitive Behavior}
\label{sec:primitive-behavior}
Some behavior cannot be expressed in the hosted language. Therefore,
such behavior has to be implemented as part of the VM. Among others,
this may be as primitive, as bytecode, or as node in an abstract
syntax tree.  For example, starting the garbage collection of the VM
the program is hosted on is behavior which is part of the VM. And
because of its low frequency of being called, it is likely implemented
as primitive.

\subsection{Speed}
The goal of reduced execution time is the main reason for creating
primitives which could be implemented in the hosted language. The
reasoning behind the reduced runtime is that every method call needs a
new frame. Even if that does not mean creating a heap object, it means
saving frame and stack pointer, copying the arguments (including the
receiver, for object oriented languages) to the new frame, and
initializing local variables. Depending on language semantics,
implementing behavior primitive may also remove costly method lookups,
repeated access boundary checks, and repeated boxing / unboxing of
basic values such as integers, floats, or strings.

Speed concerns can also be a reason to not implement behavior as
primitive. Tracing JITs record every operation in the hosted language,
but operations done by the VM in primitives may not be
visible. Implementing behavior outside of the visibility of the
optimization machinery forbids it to take those operations into
account, forcing it to create safe borders whenever leaving traces for
primitives. On traces exits, the VM has to create a valid frame for
the primitive to operate in, and, if applicable, box the arguments and
intermediate results. If a VM utilizes tracing optimizations, it needs
to be able to trace primitives. The easiest way to make a primitive
visible to the tracing machinery is not having it, i.e. implementing
it in the hosted language.

\subsection{Programming Process}
Developing a VM has different properties regarding the development
process and the ease of debugging, relative to the hosted language.
First, the round-trip times may differ. The round-rip time is measured
between editing the code and seeing changed behavior. The round-trip
time depends on features such as incremental compilation, and the
general speed of the VM. In order to see the results of a VM change,
the VM has to be recompiled, and then, the program needs to run again
to the point of failure. If the source of the hosted program changes,
the VM only has to rerun, skipping the compilation step.  During
debugging, if the VM is developed in a language which itself
translates to C, it might have the problem of translucent
abstractions. Even though programmers writes code on one abstraction
level, they need to be aware of specifics of lower abstraction
levels. And when they debug on that lower level, they might have to
match the translated sources with their own code, further increasing
the effort required, and thus impeding the ability to debug.  Last,
the steps taken within a primitive are usually hidden from the hosted
language. This has two implications. First, this means a step-wise
debugger can but step over the primitive, since it has no intermediate
representation to step through. Second, the primitives source code is
hidden in the VM sources, and usually also in a different programming
language. Thus, they require explicit documentation, and are not
easily understandable.

\subsection{Reuse}
Most widely used languages have more than one implementation. Examples
are CPython, IronPython, and PyPy, all of which implement Python, MRI,
JRuby, Rubinius, and Topaz, which implement Ruby, and Cog, RoarVM, and
R/SqueakVM, which implement Squeak Smalltalk. Even if different VMs
might not be compatible at the level of intermediate representation,
any behavior not implemented in the hosted language needs to be
supplied by the VM. Subsequently, any code which is not part of the
original VM increases the amount of work needed to reimplement that
language. In other words, any complex behavior implemented in the
hosted language based on a small set of basic primitives can be taken
and used directly, in contrast to programs which rely on hot loops
implemented primitive, or C-compatibility for VM
interfaces. Furthermore, a VM reimplementation may decide to not
adhere to C-interfaces, or to reduce the number of primitives. In such
cases, changes to the existing codebases of libraries are required.
Unfortunately, such changes break compatibility with existing language
implementations, leading to either increased effort to offer two
versions of popular libraries with only small differences for the
language implementations, or a decision to only support one of the
VMs.

- also discuss how Smalltalk fallback code duplicates VM behavior (and
may deviate)

\subsection{Changeability}
When behavior is implemented as primitive, changing it usually
requires recompiling the VM, or the according extension. Since
rebuilding requires additional sources and a build environment, this
not only affects the round- trip times, but also encumbers changing
primitives at all.  On the other hand, there might be behavior which
should not be subject to change. By fixing that behavior in form of a
primitive, it is possible to assert a certain behavior by controlling
the VM binary.

\section{Measurements}
\label{sec:measurements}

\subsection{The Primitives in Question}
\label{sec:primitives-question}

- 4

- array fill

- bitblt

- largeInteger something

\subsection{Measuring the Gap}
\label{sec:measuring-gap}

- 5.1 + 5.2

\subsection{Analysis}
\label{sec:analysis}

- 5.3 Explaining the Gap

- 5.4 Contesting the results

\section{Related Work}

- 6.1

- 6.2 + Topaz

\section{Conclusions}
\label{sec:conclusions}

- 5.5 Lessons learned

- 7 Conclusions

\bibliographystyle{abbrvnat}

% The bibliography should be embedded for final submission.

\begin{thebibliography}{}
\softraggedright

\bibitem[Smith et~al.(2009)Smith, Jones]{smith02}
P. Q. Smith, and X. Y. Jones. ...reference text...

\end{thebibliography}


\end{document}

%                       Revision History
%                       -------- -------
%  Date         Person  Ver.    Change
%  ----         ------  ----    ------

%  2013.06.29   TU      0.1--4  comments on permission/copyright notices


%%% Local Variables:
%%% mode: latex
%%% TeX-master: t
%%% End:
